\documentclass[12pt]{article}
\usepackage[utf8]{inputenc}
\usepackage[T1]{fontenc}
\usepackage{hyperref}
\title{Kayan River Kayan Language Resources for K-12 Curriculum Development}
\author{Compiled by Claude Sonnet 5.5 (claude-sonnet-5-5) for Andr\'as Kornai}
\date{}
\begin{document}
\maketitle

\section{Introduction}
Kayan River Kayan (ISO 639-3 xkn; Glottocode kaya1325) is an Austronesian language of North Kalimantan, Indonesia, with about 2,000 speakers according to the task description. It is one of several Kayan varieties (others include Kayan Mahakam xay and Busang Kayan bfg). The search found no resource specific to xkn.

\section{Primary Educational Resources}
None found.

\section{Linguistic and Reference Materials}
\begin{itemize}
\item Wikipedia on the Kayan language cluster, which lists the xkn variety: \url{https://en.wikipedia.org/wiki/Kayan_language_(Borneo)}
\item News report on Curtin Malaysia's Kayan Ethnographic Online Dictionary (May 2026). Which Kayan variety it covers is not stated, so it may not cover Kayan River Kayan: \url{https://www.theborneopost.com/2026/05/22/curtin-launches-kayan-ethnographic-online-dictionary-to-preserve-indigenous-heritage/}
\end{itemize}

\section{Religious and Community Resources}
None found.

\section{Digital and Audio Resources}
None found for this variety.

\section{Curriculum Development Context}
The language is endangered and not known to be taught in schools. Related Kayan work, such as the Curtin dictionary, may offer a model or partial vocabulary overlap. The extent of overlap with Kayan River Kayan has not been checked.
\end{document}